\section{Discussion}\label{sec:discussion} % (fold)

\subsection{Partitioned hardware}\label{sec:hardware}
While this thesis makes use of the \texttt{fence.t} instruction to clear the
microarchitectural state on a domain switch, a lot of the hardware still remains
shared. Larger caches such as the L2 and L3 cache are too large to simply
flush \cite{wistoff2023}. This means higher level caches need another method for
partitioning. In \cite{ge2019} they make use of cache coloring to partition the
higher level caches. In essence cache coloring works by using some of the bits
in the physical address to encode a color of the page. For example, if using 2
bits in the physical address, one can encode 4 colors, and when the application
asks to allocate consecutive pages, then the OS makes sure to map them to the
correct colors \cite{dinh, WhatCacheColoring}. The L1 cache is usually indexed
directly with virtual addresses, and as such can not be colored in this manner,
whereas higher-level caches more often make use of elements of the physical
address.

Cache coloring utilises the cache in an unintended manner. In general, it
reduces the average performance of the entire system, while for some specific
applications it can improve performance by reducing cache interference.
However, in truly secure kernel software partitioning is still required and as
such it is a cost that is inevitable. The larger issue with cache coloring is
the uncertainty that it poses. Where the Instruction Set Architecture(ISA)
usually standardizes the architectural interface for virtual memory and the MMU
behaviour, the same cannot be said for the underlying foundation of cache
coloring. For example, in RISC-V the ISA specifies the programmer-visible MMU
model, while the underlying cache implementation remains undefined. This means
that the cache coloring implementation is CPU-specific based on things like
cache indexing. This makes a cross platform implementation of cache coloring
challenging, and any future development in cache architecture could break any
such implementation \cite{WhatCacheColoring}.

In recent years different architectures have added support for hardware
cache partitioning. For example, Intel have added
CAT \cite{IntroductionCacheAllocation}, ARM with the MPAM
extension \cite{LearnarchitectureMemory} and similar work for RISC-V with the
QoS register interfaces \cite{11IntroductionRISCV}. For a truly secure kernel,
these items need to reach a maturity level, where domains can be guaranteed
separate in the higher level caches that previously could not be partitioned.

Another approach for safe OS execution is what was seen in \cite{karlsson2025}.
Here they make use of the Scratch Pad Memory (SPM) which is a section of memory
that is never backed up to the higher level caches. The kernel itself was then
developed, such that all memory accesses by the kernel only left a footprint in
the SPM, which was then flushed on a domain switch. This approach seems
promising for the kernel data itself, assuming that the SPM is guaranteed large
enough to hold the kernel footprint. Even when cache partitioning fully
matures, this might still be a part of the solution, as it most likely would
have performance improvements over cache partitioning, and would reduce the
interference the OS would have on the user level applications using the higher
level caches.

Most likely there will never be a hardware solution which fits all. A general
purpose operating system, where user-level constant-time programming can solve
the most important issues, might not need specific domain separation and cache
partitioning. On the other hand, in high security environments this is a must
and something where depending on the use case the correct solution could
vary.

\subsection{Reflection}
The underlying implementation is functional and via the evaluation seems to
provide the correct time slice based scheduling. However, the implementation
adds a significant amount of complexity to the FreeRTOS kernel, and this new
domain-level scheduler seems far distant from the general FreeRTOS scheduler.
The domain-level scheduler also changes some fundamental items of the priority
based scheduler. Take for example the \texttt{vTaskDelayUntil} function from
FreeRTOS. This function takes the number of ticks the current task should wait
until it can be scheduled again. This works via the global \texttt{xTicks},
which is now domain specific in the domain scheduler. This means that if a
task attempts to sleep for 2 ticks, those two ticks are only counted when the
tasks own domain is running. This breaks a lot of assumptions one might have
coming from the standard FreeRTOS implementation, and changes how some of the
fundamental timing based items function within FreeRTOS. 

So while time protection seems to be a reachable goal, it comes with significant
drawbacks. It adds to the complexity of the kernel, hinders the average
performance and restricts how the kernel operates significantly. As such, it
seems unreasonable to aim for time protection across all future kernels, and a
general operating system should most likely not strive towards full time
protection. However, some aspects of the work done could translate. As an
example, one might want to create a single task for handling encryption, where
the microarchitectural state is cleared before scheduling the next task.
Seeing that such an operation would on average only add 1\% to the scheduling,
as seen with the temporal fence instruction, this could be worth it for the
extra security. Together with user level mitigations most items deemed
sensitive could be protected.

However, one could still find situations where this is simply not enough. In
these high security environments a truly time protected kernel is still needed.
It gives guarantees that possible covert channels are not possible and in turn
guarantees the closure of timing channels. In these environments formal
verifications are also necessary. At the time of writing there is no formally verified
kernel, which guarantees time protection, which these environments would
require. 

% section Discussion (end)
